\section{Introduction}
\label{sec:intro}

Large language models answer fluently, but what they know is fixed at training
time. They have not seen documents written after their cutoff or anything
kept inside an organisation, and when a question falls outside what they
learned they tend to produce a confident wrong
answer~\cite{ji2023hallucination}. Retrieval-augmented generation
(RAG)~\cite{lewis2020rag,gao2023ragsurvey} fetches relevant passages at query
time and asks the model to answer from them, which moves the burden onto the
retriever: if the passage that holds the answer is not retrieved, the
generator cannot recover it.

The two common retriever families fail in different places. Dense
retrievers~\cite{karpukhin2020dpr} match meaning and handle paraphrase well,
but they often miss exact strings that were rare in their training data,
such as version numbers or header names. A question about a status code can
be answered from one line of one specification, and to an embedding model
that line may look much like dozens of other passages on the same topic. BM25~\cite{robertson2009bm25} finds such
strings easily, but it fails when the question and the source use different
words. Hybrid retrieval runs both and merges the results, usually with
Reciprocal Rank Fusion (RRF)~\cite{cormack2009rrf}, with one constant
balance between the two sides. That constant is a compromise between queries
that need exact matching and queries that need semantic matching.

Choosing the balance per query is not a new idea. DAT~\cite{hsu2025dat} asks
an LLM to judge the top result of each retriever and sets the weight from
those judgements. That works, but it puts LLM calls on the retrieval path of
every query. We ask whether cheap signals from the query text and the two
candidate lists are enough. Our
method, Query-Adaptive Fusion (QAF), feeds 12 such features to a ridge
regression that outputs the fusion weight. It is trained once on development
queries, needs no LLM at query time, and its coefficients can be read off
directly. In our experiments the answer is mostly no. QAF is cheap and
readable, but it does not reliably beat a well-chosen fixed weight, and we
report this result together with an analysis of why.

We build a complete pipeline, from structural
chunking through BM25 and dense indexing, weighted RRF and cross-encoder
reranking to a local 7B generator. Everything runs on a single laptop. We test
it on three public collections and on a new corpus of 20 IETF RFCs, with two
encoders: BGE-small and Google's recent EmbeddingGemma~2. Our contributions
are:
\begin{itemize}
  \item \textbf{QAF}, an LLM-free, interpretable predictor of the per-query
  dense/sparse fusion weight, trained on a plateau-centre regression
  target derived from per-query weight sweeps, together with an analysis of
  why it does not reliably beat a tuned fixed weight even though per-query
  headroom exists.
  \item \textbf{An identifier-heavy RFC test collection} of 2{,}973 chunks and
  101 questions, generated with a local LLM, filtered automatically and
  checked by all three authors, with relevance labels derived from answer
  spans.
  % NOTE(results, part 2): add "an ablation" / "EmbeddingGemma 2 at four
  % embedding sizes" back to this bullet only if Task 14 produces them.
  \item \textbf{A controlled comparison} of sparse, dense, fixed-weight hybrid
  and QAF retrieval, with and without reranking, on four collections and two
  encoders, with paired bootstrap tests. It shows that fusion helps when the
  two retrievers are comparably strong, that BM25 is the strongest single
  retriever on identifier-heavy RFC text, and that reranking can hurt a
  strong first stage.
  % NOTE(results, part 2): keep "answer quality ... judged faithfulness" only if
  % Task 15 (generation, RAGAS) produces results.
  \item \textbf{An end-to-end account of cost}: per-stage latency and answer
  quality (exact match, F1 and judged faithfulness) on commodity hardware
  with no paid API.
\end{itemize}

Section~\ref{sec:related} reviews related work. Section~\ref{sec:method}
describes the pipeline and QAF, and Section~\ref{sec:setup} the data,
configurations and metrics. Section~\ref{sec:results} reports results,
Section~\ref{sec:discussion} discusses them along with the limitations of the
study, and Section~\ref{sec:conclusion} concludes.
